\section{Design}

\subsection{Architectural sketch}

Below is the architectural sketch of the components structured in the diagram that highlights the
software and hardware components.

\begin{figure}[h]
    \centering
    \includegraphics[width=1\textwidth]{./img/architecturediagram.png}
    \caption{Architectural sketch of the project}
    \label{fig:archi}
\end{figure}

This is the description of each of the elements of the diagram:
\begin{enumerate}
  \item \textbf{ESP32 DevKit microcontroller} -- central processing unit board that runs FreeRTOS
    and executes all application tasks preemptively
  \item \textbf{Button (GPIO\,15)} -- input peripheral with internal pull-up resistor, used by Task 1
    to detect press/release transitions
  \item \textbf{Green LED (GPIO\,13)} -- output indicator for short button presses ($< 500$ ms)
  \item \textbf{Red LED (GPIO\,26)} -- output indicator for long button presses ($\geq 500$ ms)
  \item \textbf{Yellow LED (GPIO\,32)} -- output indicator that blinks rapidly to acknowledge each press
  \item \textbf{Scheduler module} -- software service that creates FreeRTOS tasks from plain task
    function descriptors, applying period and offset via \texttt{vTaskDelayUntil} and \texttt{vTaskDelay}
  \item \textbf{Sync module} -- software service that wraps a FreeRTOS binary semaphore (event
    signaling) and a mutex (shared variable protection) behind a simple API
  \item \textbf{Serial module} -- software driver that initializes UART at 115200 baud; on ESP32,
    \texttt{printf()} is natively routed to the serial interface
  \item \textbf{LED driver module} -- software driver providing abstractions for controlling each LED
    independently
  \item \textbf{Button driver module} -- software driver providing button state reading
\end{enumerate}

\subsection{Electrical sketch}

The circuit diagram outlines the setup for this project, showing the connections between
the ESP32, three LEDs (green, red, yellow), a push button, and the computer for serial
communication. Each LED is connected to its respective GPIO pin through a $220\,\Omega$
current-limiting resistor. The push button is connected to GPIO\,15 with the internal pull-up
resistor enabled, meaning the pin reads LOW when the button is pressed.

\begin{figure}[h]
    \centering
    \includegraphics[width=0.4\textwidth]{./img/circuit_sketch_wokwi.jpg}
    \caption{Circuit sketched on wokwi}
    \label{fig:wokwi}
\end{figure}

Figure \ref{fig:wokwi} details the connections and components that make up the system. This visual
documentation is crucial for replicating the setup or diagnosing issues, as it clearly shows
how each component is integrated into the overall design.

\subsection{Flow chart}

The Figure \ref{fig:flow} illustrates the FreeRTOS preemptive scheduling flow. The RTOS kernel
manages three independent tasks, each running on its own period. Task 1 polls the button every
10\,ms and, on a press release, writes shared data under mutex and gives the binary semaphore.
Task 2 runs every 50\,ms and non-blockingly checks the semaphore; if taken, it reads press data
and updates statistics under mutex, then drives the yellow LED. Task 3 runs every 10\,seconds,
acquires the mutex to snapshot and reset statistics, and prints the report.

\begin{figure}[h]
    \centering
    \includegraphics[width=1\textwidth]{./img/flow.png}
    \caption{Flow chart of the algorithm}
    \label{fig:flow}
\end{figure}

\clearpage
\subsection{Code design}

Code is developed following the principle of separating the interface from the
implementation, by creating for each header file \texttt{.h} a respective
\texttt{.cpp} file that implements it. The project is organized into three layers:

\begin{enumerate}
  \item \textbf{Driver layer} -- contains hardware-specific code for the button
    (\texttt{button.h/cpp}) and LEDs (\texttt{led.h/cpp}). All direct \texttt{Arduino.h} calls
    are encapsulated here.
  \item \textbf{Service layer} -- contains two reusable, hardware-independent modules:
    the scheduler (\texttt{scheduler.h/cpp}) which creates and manages FreeRTOS tasks, and
    the sync service (\texttt{sync.h/cpp}) which encapsulates the binary semaphore and mutex.
    The serial service (\texttt{serial.h/cpp}) initializes UART.
  \item \textbf{Application layer} -- contains the task logic (\texttt{lab4.h/cpp}) which defines
    the three plain one-shot task functions, their period/offset/priority parameters, shared global
    variables, and the setup routine.
\end{enumerate}

The scheduler uses a \texttt{TaskDef} structure array passed to \texttt{srvSchedulerSetup()},
which creates a FreeRTOS task for each entry:
\begin{lstlisting}
typedef struct {
  const char      *name;
  TaskFunc         func;       // plain void() function -- same shape as bare-metal
  uint32_t         periodMs;   // recurrence period
  uint32_t         offsetMs;   // initial delay before first run
  uint32_t         stackSize;
  UBaseType_t      priority;
} TaskDef;
\end{lstlisting}

Internally, the scheduler wraps each plain function in a \texttt{periodicTaskRunner} that
applies the offset via \texttt{vTaskDelay} and then loops with \texttt{vTaskDelayUntil}
for drift-corrected periodic execution. This design keeps the infinite loop entirely inside
the scheduler abstraction, not inside the task functions.

Tasks communicate and synchronize through the sync service abstraction:
\begin{lstlisting}
// Binary semaphore -- event signaling (Task 1 -> Task 2)
void syncSemaphoreGive();
bool syncSemaphoreTake(uint32_t timeoutMs);  // timeoutMs=0 is non-blocking

// Mutex -- shared variable protection (Tasks 1, 2, 3)
void syncMutexTake();
void syncMutexGive();
\end{lstlisting}

Task 1 calls \texttt{syncMutexTake/Give} to write \texttt{lastPressDuration} and
\texttt{lastPressWasShort}, then calls \texttt{syncSemaphoreGive} to notify Task 2.
Task 2 calls \texttt{syncSemaphoreTake(0)} (non-blocking, timeout\,=\,0) to check for a
pending event, then uses the mutex to read press data and update counters.
Task 3 uses the mutex to atomically snapshot and reset all statistics counters before printing,
ensuring no partial update from Task 2 is observed.

This design ensures that \texttt{printf()} is called from within the Task 3 FreeRTOS task
context, never from an ISR, avoiding blocking I/O in interrupt context.
